\section{モデリングの目から見た検定2;パラメータの世界とデータの世界}
\label{lesson2_26}
\subsection{授業内容}
\subsubsection{科目の中でのこのコマの位置づけ}
ここまでGLMM，HLMとさまざまな分布を組み合わせて利用するモデルについてみてきた。

今回は混合正規分布モデルと0過剰ポアソンモデルを考える。これまではデータが全体的に均質的であることを想定していたが，そもそも異なる種類のデータが混じり合っていると考えられる場合，異なる分布をあてがう方が良い。
ここでどちらの分布に属するかがある種の確率で代わり，それに従って続くモデルが異なるという，離散確率分布による条件分岐をデータ生成メカニズムに導入することになる。
またこれをStanで実行する場合には，離散変数が直接扱えないことから，すべての可能な場合を数え上げて足し合わせるという周辺化して消去するというテクニックを利用することになる。技術的に高度な側面もあるが，条件分岐をデータ生成メカニズムに組み合わせることができれば表現力は一気に広がることになる。
\subsubsection{コマ主題細目}
\begin{description}
	\item[混合分布モデル] 混合分布モデルは確率的なクラスター分析(Model Based Clustering)でもある。階層線形モデルと違って，クラスター分析はデータの背後にあるクラスが明示的に示されておらず，データの適合から考えることになる。まずデータの可視化によってまずその可能性に気づき，これをどのようにモデル化できるか，そのアイデアを理解する。具体的には，ベルヌーイ分布など離散変数のパラメータによって条件分岐が発生し，各条件のもとで確率モデルが描かれることになるが，この確率モデルをどのように統合するかということを設計図の段階で理解する。設計図での理解を踏まえて，Stanでの実装レベルでの理解に進むことができるのだから。
	      \begin{flushright}
		      $\to$\textcite{Matsuura201610}のPp.209--213.
	      \end{flushright}

	\item[周辺化消去] Stanは離散確率分布を直接モデルの中に組み込むことはできない。そこで\verb|log_sum_exp|という特殊な関数を使って，すべての場合わけを行った離散モデルの統合を考える。まずターゲット記法について理解し，続いて周辺化消去の書き方を理解する。また混合正規分布モデルの場合，ラベルスイッチングの問題が生じることが考えられるから，\verb|ordered vector|の型を利用することが多い。こうした特殊な関数やベクトルについても理解を深める
	      \begin{flushright}
		      $\to$\textcite{Matsuura201610}のPp.203--208.
	      \end{flushright}
	\item[0過剰ポアソン] データの性質を考えると，必ずしも正規分布モデルばかりではなく，離散変数など一般的なモデルまで拡張することが可能である。そこで野球データなど具体的な分布情報とともに，ゼロ過剰ポアソン分布を使ったモデルを考える。データ生成のメカニズムがより具体的，実践的に考えることができるので，モデリングによる分析の自由度が高まることを実感できるだろう。
	      \begin{flushright}
		      $\to$\textcite{Matsuura201610}のPp.82，83，85--89.
	      \end{flushright}
\end{description}
\subsubsection{キーワード}
\begin{itemize}
	\item クラスター分析
	\item 混合分布モデル
	\item 周辺化消去
	\item ゼロ過剰ポアソン分布
\end{itemize}
\subsection{授業情報}
\paragraph{コマの展開方法}
講義/遠隔可/演習
\subsubsection{予習・復習課題}
\paragraph{予習}データの描画から気づかされることは無数にある。探索的にデータプロットができるように，\verb|ggplot|やデータハンドリング技術について復習しておくことが望ましい。
\paragraph{復習}混合分布モデルが応用できるようなデータ例を身の回りから考えてみるとよい。その上で，データがどのような背景から生成されているかの設計図を書き，設計図からコードに起こすという手順を一歩ずつ確認しておこう。